\chapter{Methodology}
% ═══════════════════════════════════════════════════════════════════════════════



% edit below this
\subsection{BB84}
\label{sec:bb84}
\begin{figure}[ht]
    \centering
    \includegraphics[width=0.6\linewidth]{images.methodology/QuantumAbstractTNR.png}
    \caption{Abstract P\&M QKD Network Model}
    \label{fig:QKDnetworkmodel}
\end{figure}

The BB84 protocol, proposed by Bennett and Brassard in 1984 \cite{BB84}, was the first practical QKD scheme. It is a prepare-and-measure protocol that encodes classical bit values onto single-photon polarisation states drawn from two mutually unbiased bases: the rectilinear basis ($\ket{0},\ket{1}$) and the diagonal basis ($\ket{+},\ket{-}$). The sender, Alice, randomly selects both a bit value and a basis for each qubit, producing one of four possible states:
%
\begin{equation}
  \ket{0},\quad \ket{1},\quad
  \ket{+} = \frac{\ket{0}+\ket{1}}{\sqrt{2}},\quad
  \ket{-} = \frac{\ket{0}-\ket{1}}{\sqrt{2}}.
  \label{eq:bb84-states}
\end{equation}
%


These states are transmitted to Bob over the quantum channel. Bob then independently and randomly chooses a measurement basis for each incoming qubit. Because the two bases are mutually unbiased, a measurement in the wrong basis yields a uniformly random outcome and provides no information about Alice's bit value \cite{QKDbible}.

After transmission is complete, Alice and Bob switch to an authenticated classical channel to perform \textit{sifting}, as seen in Fig. \ref{fig:QKDnetworkmodel}. At this point, Bob announces only his choice of basis for each qubit, and Alice confirms which choices matched her preparation basis. All rounds in which the bases disagree are discarded. In an ideal, noiseless channel, the remaining \textit{sifted key} is perfectly correlated between Alice and Bob, and roughly half the transmitted qubits survive the sifting step \cite{BB84}.

\begin{figure}[h]
    \centering
    \includegraphics[width=0.7\linewidth]{images.methodology/bb84processTNR.png}
    \caption{BB84 Operation Model}
    \label{fig:BB84operation}
\end{figure}

To check for eavesdropping, Alice and Bob sacrifice a random subset of the sifted key and compare bit values over the public channel. Any intercept-resend attack by an eavesdropper, Eve, necessarily disturbs a fraction of the qubits because she cannot simultaneously determine the polarisation in both bases without violating the no-cloning theorem \cite{QKDbible}. A full intercept-resend strategy is expected to produce a Quantum Bit Error Rate (QBER) of approximately 25\%, well above the theoretical security threshold of approximately 11\% \cite{ekertQI}. If the observed QBER exceeds this threshold, the protocol aborts. Otherwise, Alice and Bob proceed through error correction and privacy amplification to yield a final secret key, depicted in Fig. \ref{fig:BB84operation}.


\subsection{B92}
\label{sec:b92}

\begin{figure}[h]
    \centering
    \includegraphics[width=0.6\linewidth]{images.methodology/b92basis.png}
    \caption{B92 Measurements Bases}
    \label{fig:B92bases}
\end{figure}

In 1992, Bennett proposed the B92 protocol as a simplification of BB84 \cite{B92}. The central insight is that unconditional security does not require four encoding states; two non-orthogonal states (seen Fig. \ref{fig:B92bases} are sufficient because no measurement can perfectly distinguish between non-orthogonal quantum states \cite{ekertQI}. In B92, Alice encodes bit value 0 as $\ket{0}$ and bit value 1 as $\ket{+}$:

%
\begin{equation}
  \ket{\varphi} =
  \begin{cases}
    \ket{0}  & \text{if } a = 0, \\[4pt]
    \ket{+}  & \text{if } a = 1.
  \end{cases}
  \label{eq:b92-states}
\end{equation}
%

Bob generates his own random bit $a'$ and measures in the $Z$ basis ($\ket{0},\ket{1}$) if $a' = 0$, or in the $X$ basis ($\ket{+},\ket{-}$) if $a' = 1$. When Bob's measurement basis is incompatible with Alice's state, the outcome is conclusive only some of the time, hence, inconclusive results (a.k.a. erasures) are discarded \cite{QKDbible}. Bob publicly announces which rounds yielded a conclusive result, and Alice retains only the corresponding bits. Neither party reveals their basis choice or bit value during this exchange.

Because B92 relies on only two states rather than four, its sifting efficiency is lower than that of BB84, and its tolerance to channel noise is reduced. The theoretical QBER security threshold for B92 is approximately 6.5\%, compared to approximately 11\% for BB84 \cite{ekertQI}. Furthermore, the two-state encoding makes B92 more susceptible to unambiguous state discrimination (USD) attacks, in which Eve exploits the known non-orthogonal structure to occasionally identify Alice's state with certainty while flagging ambiguous cases as losses \cite{B92}. Despite these limitations, B92 remains an important protocol both for its theoretical elegance and because it demonstrates that the minimal resource needed for secure QKD is a pair of non-orthogonal states.

\subsection{E91}
\label{sec:e91}

In contrast to the P\&M paradigm of BB84 and B92, the E91 protocol proposed by Ekert in 1991 \cite{E91} derives its security from quantum entanglement and Bell's theorem. A source prepares pairs of qubits in the singlet Bell state:

%
\begin{equation}
  \ket{\Psi^{-}} = \frac{\ket{01} - \ket{10}}{\sqrt{2}},
  \label{eq:e91-singlet}
\end{equation}
%

and distributes one qubit from each pair to Alice and the other to Bob. Each party independently selects a random measurement orientation from a predetermined set of angles. When Alice and Bob happen to choose the same measurement basis, their outcomes are perfectly anti-correlated, and these matching rounds are used to construct the raw key \cite{E91}.

The rounds in which Alice and Bob chose different bases serve a separate purpose: they are used to evaluate the Clauser-Horne-Shimony-Holt (CHSH) inequality \cite{QKDbible}. For the singlet state, quantum mechanics predicts a CHSH parameter of $S = 2\sqrt{2}$, which maximally violates the classical bound of $S \leq 2$. Any eavesdropping that disturbs the entanglement between the pairs will reduce the observed value of $S$ towards or below the classical limit, thereby revealing Eve's presence \cite{ekertQI, E91}. This connection between eavesdropper detection and Bell inequality violation is the defining feature of E91 and distinguishes it fundamentally from P\&M protocols.

\begin{table}[h]
\centering
\caption{Qualitative Comparison of BB84, B92, and E91}
\label{tab:protocol-comparison}
\footnotesize
\begin{tabularx}{\linewidth}{@{}l X X X@{}}
\toprule
\textbf{Feature} & \textbf{BB84} & \textbf{B92} & \textbf{E91} \\
\midrule
Paradigm & Prepare \& Measure & Prepare \& Measure & Entanglement-Based \\
Encoding states & 4 (two bases) & 2 non-orthogonal & Bell pairs, 3 angles \\
Alice's bases / angles & $\{Z, X\}$ & $\{\ket{0}, \ket{+}\}$ & $\{0, \pi/4, \pi/2\}$ \\
Bob's bases / angles & $\{Z, X\}$ & $\{Z, X\}$ & $\{\pi/4, \pi/2, 3\pi/4\}$ \\
Sifting efficiency & $\sim$50\% & $\sim$25\% & $\sim$22.2\% \\
Threshold ($f_{EC}\!=\!1.0$) & $\sim$11\% & $\sim$8\% & $\sim$11\% \\
Threshold ($f_{EC}\!=\!1.16$) & $\sim$9.8\% & $\sim$6.5\% & $\sim$9.8\% \\
Eve detection & QBER & QBER & CHSH $+$ QBER \\
Full intercept-resend QBER & $\sim$25\% & $\sim$33--38\% & degraded $|S|$ \\
Known vulnerabilities & PNS (WCP sources) & USD attack & Entanglement loss \\
Hardware complexity & Moderate & Lowest & Highest \\
\bottomrule
\end{tabularx}
\end{table}

This entanglement-based approach offers a conceptual advantage as the key bits are not determined until the moment of measurement, so there is no point during transmission at which an encoded classical value exists for Eve to intercept. However, E91 places greater demands on hardware because it requires a source of high-fidelity entangled pairs and the ability to distribute them over the quantum channel while preserving entanglement \cite{ekertQI}. In a simulation context, this translates to the additional complexity of modelling entangled state preparation, Bell-state measurement, and CHSH parameter estimation alongside the standard post-processing pipeline shared with BB84 and B92.



% discuss the basis sifitng & priavcy amplification processes
\section{The Post-Processing Pipeline}
Once the quantum transmission phase of a QKD protocol is complete, the raw measurement outcomes held by Alice and Bob must be transformed into a shared secret key through a sequence of classical post-processing steps. This pipeline is common to all three protocols considered in this project, although the details of sifting differ between prepare-and-measure and entanglement-based schemes.

The first stage is \textit{sifting}, or basis reconciliation. In BB84, Alice and Bob publicly compare their basis choices and discard all rounds where their bases did not match, retaining approximately 50\% of transmitted qubits \cite{BB84}. However, in B92, Bob announces which rounds yielded a conclusive measurement outcome, and only those rounds are kept, producing a sifting efficiency of approximately 25\% \cite{B92}. For E91, sifting separates the rounds into two groups based on whether Alice and Bob selected the same measurement angle. Matching-angle rounds contribute to the raw key, while non-matching rounds are reserved for CHSH parameter estimation \cite{E91}.

Following sifting, Alice and Bob perform \textit{error estimation} by publicly comparing a randomly chosen subset of the sifted key. The fraction of disagreements in this subset provides an estimate of the QBER. If the observed QBER exceeds the protocol-specific security threshold, the key is deemed compromised and the protocol aborts. Otherwise, the sacrificed bits are discarded and the remaining sifted key proceeds to error correction \cite{Scarani2009}.

\textit{Error correction} reconciles the remaining discrepancies between Alice's and Bob's key strings. The CASCADE protocol, introduced by Brassard and Salvail \cite{Cascade1994}, is the most widely used scheme in QKD implementations. CASCADE operates through multiple passes of binary search, where in each pass, the key is divided into blocks of increasing size, and parity checks identify and correct single-bit errors within each block. The multi-pass structure progressively eliminates errors that were missed in earlier rounds. A critical parameter is the error correction efficiency $f_{EC}$, defined as the ratio of information actually leaked during correction to the Shannon-limit minimum. CASCADE achieves $f_{EC} \approx 1.16$ in practice \cite{Cascade1994}, meaning it leaks approximately 16\% more information than the theoretical minimum. This distinction has direct consequences for the achievable secure key rate, as discussed below.

The final stage is \textit{privacy amplification}, which compresses the corrected key to remove any information that may have leaked to an eavesdropper during both the quantum transmission and the error correction phases. Universal hashing functions are applied to map the corrected key of length $n$ to a shorter final key of length $\ell < n$, where $\ell$ is chosen such that Eve's mutual information with the final key is negligibly small \cite{Bennett1995PA}. The amount of compression required depends on both the estimated QBER and the information leaked during error correction.

The GLLP formula \cite{GLLP2004} quantifies how much key material survives this entire pipeline. For BB84 with a single-photon source, the asymptotic secure key rate per sifted bit is given by
%
\begin{equation}
\label{eq:gllp}
R = 1 - h(\text{QBER}) - f_{EC} \cdot h(\text{QBER}),
\end{equation}
%
where $h(\cdot)$ denotes the binary Shannon entropy. Setting $f_{EC} = 1.0$ (the Shannon limit) yields the well-known 11\% QBER threshold for BB84. However, using the realistic CASCADE efficiency of $f_{EC} = 1.16$ shifts this threshold to approximately 9.8\%, a reduction with meaningful engineering consequences for budget calculations. It is worth noting that many simulation studies in the literature default to $f_{EC} = 1.0$, which overestimates the achievable secure key rate. This project uses $f_{EC} = 1.16$ throughout to provide a more accurate representation of practical system performance.

\section{Noise Models for Modelling} 

In any practical QKD deployment, the quantum channel introduces errors that are physically indistinguishable from those caused by eavesdropping. Accurately modelling these channel imperfections is therefore essential for a simulation platform that seeks to evaluate protocol performance under realistic conditions.

The \textit{depolarising channel} is the most commonly used noise model in QKD analysis. It maps an arbitrary single-qubit state $\rho$ to the maximally mixed state with probability $p$, effectively applying a random Pauli error:
%
\begin{equation}
\label{eq:depolarising}
\mathcal{E}_{dep}(\rho) = (1-p)\rho + \frac{p}{3}(X\rho X + Y\rho Y + Z\rho Z).
\end{equation}
%
Because the depolarising channel applies errors isotropically across all axes of the Bloch sphere, it affects all encoding bases equally, Physically, it represents a worst-case model of isotropic noise arising from environmental decoherence in the fibre channel.

\section{QKD Simulation Frameworks}
Qiskit Aer \cite{QiskitAer}, the noise simulation back-end of IBM's Qiskit framework, provides gate-level circuit simulation with a built-in noise model library, compatible with Python3. Its \texttt{NoiseModel} API allows noise channels to be attached to specific gates rather than applied globally, enabling the channel-marker approach used in this project: an identity gate is inserted between state preparation and measurement to represent the quantum channel. By applying the noise exclusively to this gate, the trusted-device assumption that Alice's preparation and Bob's measurement remain noise-free is preserved. 

\section{Research Gaps}





% >>> ADDED SECTION: Simulation Architecture <<<
\section{Simulation Architecture}

% discuss the process of developing the system thru AGILE methodologies, starting priorities lied within correctly implementing protocol logic as opposed to focusing on large system features

\subsection{Modular Package Design}

The simulation platform is implemented as a modular Python package (\texttt{qkd\_sim}), structured across nine source files with clearly separated uses; protocol logic, noise model configuration, eavesdropper simulation, benchmarking orchestration, visualisation, and a command-line interface (CLI). This modular architecture was a deliberate design decision driven by two requirements. First, protocol-agnostic benchmarking: by separating the noise and eavesdropper modules from individual protocol implementations, the same channel conditions can be applied uniformly across BB84, B92, and E91 without code duplication. Second, reproducibility: decoupling the simulation parameters from the execution code ensures that any experiment can be fully reproduced from its configuration file alone.

The system employs OOP, organised around an abstract base class, \texttt{QKDProtocol}, which defines the common interface shared by all protocol implementations. Each protocol subclass (BB84, B92, E91) overrides the abstract \texttt{run()} method while inheriting the shared post-processing logic. A companion dataclass, \texttt{QKDResult}, encapsulates the outputs of every protocol run, including Alice's and Bob's sifted keys, the computed QBER, and sifting efficiency. This design ensures that the benchmarking module can treat all protocols interchangeably, enabling fair comparison under identical conditions.

% NOTE: Replace with your actual source listing or figure
\begin{figure}[h]
    \centering
    \includegraphics[width=0.6\linewidth]{images.methodology/system overview.png}
    \caption{Module dependency structure of the system.}
    \label{fig:architecture}
\end{figure}

\subsection{Configuration System}
Initially, simulation parameters were hardcoded into the system to accelerate the process of implementing the protocol logic. Through AGILE development methodologies, I was able to expand the capability of the system by implementing configurable YAML files to load simulation parameters into the system. YAML was selected over JSON for its support of inline comments, which allows each parameter to be annotated with its physical meaning and valid range. The configuration is loaded into a \texttt{SimConfig} dataclass that validates the execution mode against the four supported modes: \texttt{single} (one protocol trial with fixed parameters), \texttt{sweep} (noise strength sweep), \texttt{eve\_sweep} (eavesdropper interception rate sweep), and \texttt{protocol\_comparison} (identical noise sweep applied to multiple protocols simultaneously). A representative configuration snippet is shown in \cref{lst:yaml-config}.

% NOTE: Replace with \lstinputlisting pointing to an actual config file
\begin{lstlisting}[language=Python, caption={Example YAML configuration for a BB84 noise sweep}, label={lst:yaml-config}]
# Placeholder: insert actual YAML config snippet
# mode: sweep
# protocol: bb84
# n_qubits: 1000
# n_trials: 50
# noise_type: depolarizing
# noise_min: 0.0
# noise_max: 0.30
# noise_steps: 21
# f_ec: 1.16
# e91_channel_topology: both
\end{lstlisting}

This configuration-driven approach directly supports objective O5.3 (reproducibility) by ensuring that every simulation run is fully parameterised and traceable.

\subsection{Execution Pipeline}

The simulation follows a linear execution pipeline, illustrated in \cref{fig:pipeline}: a YAML configuration file is parsed by the \texttt{SimConfig} loader to select the protocol, channel parameters, and execution mode; the CLI dispatcher in \texttt{\_\_main\_\_.py} delegates to the appropriate runner function; Alice prepares quantum states according to the selected protocol; an identity gate is inserted as a channel marker to which noise models are attached; the optional eavesdropper module intercepts a configurable fraction of qubits; Bob performs his measurement; classical sifting extracts the raw key; and the post-processing logic computes the QBER, applies the GLLP key rate formula, and returns a \texttt{QKDResult} object. For sweep and comparison modes, the \texttt{BenchmarkRunner} orchestrates repeated trials across parameter ranges and aggregates the results into \texttt{BenchmarkData} objects containing mean and standard deviation arrays, which are then passed to the \texttt{QKDPlotter} for visualisation.

\begin{figure}[ht]
    \centering
    \includegraphics[width=0.6\linewidth]{images.methodology/information flow.png}
    \caption{Simulation execution pipeline. Noise is applied exclusively to the channel stage via the identity gate marker, preserving the trusted-device assumption.}
    \label{fig:pipeline}
\end{figure}
% >>> END ADDED SECTION <<<


% >>> ADDED SECTION: Protocol Implementations <<<
\section{Protocol Implementations}
\label{sec:protocol-implementations}

This section describes the implementation of each QKD protocol within the simulation platform. All three protocols share the same post-processing pipeline and noise application mechanism; the differences lie in state preparation, measurement, and sifting logic.

\subsection{BB84 Protocol}

The BB84 implementation follows the four-state, two-basis scheme described in \cref{sec:bb84}. For each of $n$ qubits, Alice randomly selects a bit value $a_i \in \{0,1\}$ and a basis $\hat{a}_i \in \{Z, X\}$ using NumPy's pseudorandom number generator, then prepares the corresponding quantum state on a single-qubit Qiskit circuit. If the selected bit is 1, a Pauli-X gate is applied to flip the qubit from $\ket{0}$ to $\ket{1}$. If the diagonal basis is selected, a Hadamard gate is subsequently applied to rotate into the $X$ basis. An identity gate (\texttt{qc.id(0)}) is appended after state preparation to serve as the channel marker, to which noise models are attached during execution. Bob independently selects a random measurement basis $\hat{b}_i$; if he chooses the diagonal basis, a Hadamard gate is applied before the computational-basis measurement. All $n$ circuits are executed in a single batch call to the Qiskit Aer backend with \texttt{shots=1} and \texttt{memory=True} to retrieve individual outcomes.

The sifting procedure retains only the rounds where $\hat{a}_i = \hat{b}_i$, implemented as a Boolean mask comparison of the two basis arrays. On an ideal channel, this produces a sifting efficiency of approximately 50\% and a QBER of 0\%, serving as the primary validation targets for this protocol. The key rate is computed as the ratio of sifted key length to total transmitted qubits.

% NOTE: Replace firstline/lastline with actual line numbers from bb84.py
\lstinputlisting[
    language=Python,
    caption={BB84 state preparation and channel-marker insertion (\texttt{bb84.py}).},
    label={lst:bb84-prep},
    firstline=1,
    lastline=1
]  {}% TODO: set correct path + line range

The high-level flow of the BB84 key sifting procedure is presented in \cref{alg:bb84}.

\begin{algorithm}[H]
    \caption{BB84 Key Sifting Procedure}\label{alg:bb84}
    \KwIn{Alice's bits $a_i$, bases $\hat{a}_i$; Bob's measurements $b_i$, bases $\hat{b}_i$}
    \KwOut{Sifted key $K$}
    $K \leftarrow \emptyset$\;
    \For{$i \leftarrow 1$ \KwTo $n$}{
        Alice publishes $\hat{a}_i$ over classical channel\;
        Bob publishes $\hat{b}_i$ over classical channel\;
        \If{$\hat{a}_i = \hat{b}_i$}{
            Append $b_i$ to $K$ \Comment{Matching basis -- keep bit}
        } 
    }
    \Return{$K$}
\end{algorithm}


\subsection{B92 Protocol}

The B92 implementation uses only two non-orthogonal encoding states: $\ket{0}$ for bit value 0 and $\ket{+}$ for bit value 1. Alice prepares $\ket{0}$ directly (no gate applied) or applies a Hadamard gate to produce $\ket{+}$, followed by the identity gate channel marker (\texttt{qc.id(0)}).

Bob's measurement logic is the most distinctive aspect of the B92 implementation. Bob generates a random bit $b_i \in \{0, 1\}$ and measures in the $Z$ basis if $b_i = 0$ or the $X$ basis if $b_i = 1$, applying a Hadamard gate before measurement in the latter case. An outcome is \textit{conclusive} only if, given the basis chosen, the observed result is inconsistent with one of Alice's two possible preparations. Measuring $\ket{0}$ in the $Z$ basis can only yield $0$ and is therefore ambiguous, as is measuring $\ket{+}$ in the $X$ basis (which also yields $0$ under the Hadamard-then-$Z$ implementation). The conclusive events are the \emph{complementary} pairings: a $Z$-basis outcome of $1$ rules out a $\ket{0}$ preparation and therefore identifies Alice's bit as $1$; an $X$-basis outcome of $1$ (equivalent to projecting onto $\ket{-}$) rules out $\ket{+}$ and identifies Alice's bit as $0$. Inconclusive rounds (outcome $0$) are discarded.

The sifting step in B92 therefore differs from BB84: Bob announces which rounds were conclusive, not which basis he used. Neither party reveals their bit value or basis choice during this exchange. Bob's sifted key is inferred as $1 - b_i$ (the complement of Bob's measurement basis for the conclusive rounds), while Alice retains her original bit values for those rounds. This sifting mechanism produces a theoretical sifting efficiency of approximately 25\%, half that of BB84. The reduced efficiency, combined with the two-state encoding, results in a lower noise tolerance.

\begin{figure}[ht]
    \centering
    \includegraphics[width=0.6\linewidth]{images.methodology/b92 flowchart.png}
    \caption{B92 algorithm flowchart}
    \label{fig:pipeline}
\end{figure}

% NOTE: Replace firstline/lastline with actual line numbers from b92.py
\lstinputlisting[
    language=Python,
    caption={B92 conclusive-measurement decoding and sifting (\texttt{b92.py}).},
    label={lst:b92-sift},
    firstline=1,
    lastline=1
]{}  % TODO: set correct path + line range


\subsection{E91 Protocol}

The E91 implementation models the entanglement-based protocol described in \cref{sec:e91}. For each round, a two-qubit Qiskit circuit is initialised and a singlet Bell state $\ket{\Psi^{-}} = (\ket{01} - \ket{10})/\sqrt{2}$ is prepared by applying the following gate sequence: a Hadamard gate on qubit 0, a CNOT gate with qubit 0 as control and qubit 1 as target, a Pauli-X gate on qubit 0, and a Pauli-Z gate on qubit 0. This sequence produces the singlet state rather than the more commonly used $\ket{\Phi^+}$, matching the original E91 specification \cite{E91}. Qubit 0 is designated as Alice's and qubit 1 as Bob's.

The implementation supports two channel topologies via a configurable parameter. In the \texttt{bob} topology, an identity gate channel marker is applied only to Bob's qubit (\texttt{qc.id(1)}), representing a scenario where noise acts on the distribution channel to Bob while Alice's qubit remains local. In the \texttt{both} topology, identity gates are applied to both qubits (\texttt{qc.id(0)} and \texttt{qc.id(1)}), modelling noise on both distribution paths. This distinction affects the relationship between noise strength and observed QBER: under depolarising noise with parameter $p$, the \texttt{bob} topology produces $\text{QBER} = p/2$ (matching BB84), while the \texttt{both} topology produces $\text{QBER} = p - p^2/2$.

\begin{figure}[ht]
    \centering
    \includegraphics[width=0.6\linewidth]{images.methodology/e91 bloch sphere.png}
    \caption{E91 XZ plane measurement bases}
    \label{fig:pipeline}
\end{figure}

Alice and Bob each independently select a measurement angle from predetermined sets. Alice's angles are drawn from $\{0, \pi/4, \pi/2\}$ and Bob's from $\{\pi/4, \pi/2, 3\pi/4\}$, each chosen uniformly at random per round. Measurement in an arbitrary angle $\theta$ is implemented by applying a rotation gate $R_y(-\theta)$ before a $Z$-basis measurement. A predefined $3 \times 3$ sift table, indexed by Alice's and Bob's angle choices, classifies each round as contributing to the key, to the CHSH parameter estimation, or as discarded. Of the nine possible angle-pair combinations, two contribute to the key (matching angles that yield anti-correlated outcomes), four contribute to CHSH estimation, and three are discarded, producing a theoretical sifting efficiency of $2/9 \approx 22.2\%$.

The CHSH parameter is computed as
%
\begin{equation}
\label{eq:chsh-computation}
S = +E(a_1, b_1) - E(a_1, b_3) + E(a_3, b_1) + E(a_3, b_3),
\end{equation}
%
where $E(a_i, b_j)$ is the correlation coefficient for the angle pair $(a_i, b_j)$, calculated as the expectation value of the product of Alice's and Bob's outcomes (mapped to $\pm 1$ via $1 - 2 \times \text{outcome}$) across all rounds using that pair. The implementation is validated by confirming that $|S| \rightarrow 2\sqrt{2} \approx 2.828$ in the absence of noise or eavesdropping, consistent with the maximal quantum violation of the CHSH inequality. Under depolarising noise, the theoretical CHSH value is $2\sqrt{2}(1-p)$ for the \texttt{bob} topology and $2\sqrt{2}(1-p)^2$ for the \texttt{both} topology.

% NOTE: Replace firstline/lastline with actual line numbers from e91.py
\lstinputlisting[
    language=Python,
    caption={E91 Bell-state preparation and CHSH computation (\texttt{e91.py}).},
    label={lst:e91-chsh},
    firstline=1,
    lastline=1
]{}  % TODO: set correct path + line range
% >>> END ADDED SECTION <<<


% >>> ADDED SECTION: Modelling Assumptions <<<
\section{Modelling Assumptions}
\label{sec:modelling-assumptions}

The simulation platform operates under a set of assumptions that define the scope and limitations of the model. Stating these explicitly is important because they determine which physical effects are captured and which are deliberately excluded.

The platform assumes a perfect single-photon source where each transmission round involves exactly one photon, with no multi-photon pulses or vacuum events. Real-world QKD implementations typically use weak coherent pulse sources, which emit multi-photon pulses with non-zero probability. Multi-photon events enable photon-number-splitting attacks that are not modelled here.

All noise is assumed to originate in the quantum channel. Alice's state preparation and Bob's measurement apparatus are treated as perfect systems. This is enforced in the implementation by attaching noise models exclusively to the identity gate channel marker, ensuring that preparation and measurement operations remain noise-free. This assumption is standard in the security proofs of BB84, B92, and E91 \cite{ShorPreskill2000, Scarani2009}.

The authenticated classical channel used for sifting, error estimation, and error correction is assumed to be tamper-proof and perfectly secure. While this is a standard assumption in QKD security analysis, it is worth noting that it requires a pre-shared authentication key, creating a problem that is outside the scope of this simulation.

Finally, the model does not account for detector efficiency mismatch, timing jitter, or other side-channel imperfections. These effects are relevant to device-security analysis but are excluded here to maintain focus on the protocol-level performance as the primary objective.


% >>> ADDED SECTION: Noise Models and Channel Simulation <<<
\section{Noise Models and Channel Simulation}
\label{sec:noise-impl}

This section describes how the noise models introduced in the literature review are implemented within the simulation platform.

\subsection{Channel Marker Approach}

Noise is applied to the quantum channel using Qiskit Aer's \texttt{NoiseModel} API. The \texttt{create\_backend()} function in the noise module constructs an \texttt{AerSimulator} instance with a \texttt{NoiseModel} that attaches the specified quantum error exclusively to the \texttt{'id'} gate. An identity gate (\texttt{qc.id()}) is inserted into each single-qubit circuit between Alice's state preparation and Bob's measurement. This gate performs no logical operation but serves as a marker to which noise channels are attached. When the circuit is executed on the Qiskit Aer simulator with the configured \texttt{NoiseModel}, noise is applied exclusively to this identity gate.

This approach is critical for maintaining the trusted-device assumption. If noise were applied globally to all gates, the Hadamard gates used in Alice's preparation and Bob's measurement would also be affected, conflating device imperfection with channel noise and artificially inflating the observed QBER. By isolating noise application to the channel marker, the simulation accurately represents a scenario where only the quantum transmission link is imperfect.

The noise module supports the ideal \& depolarising channel for the benchmarking analysis, as justified in the literature review.

\subsection{Eavesdropper Model}

The eavesdropper simulation is implemented in the \texttt{EveInterceptor} class, which performs an intercept-resend attack with a configurable interception probability $p \in [0, 1]$. For each qubit in transit, Eve intercepts it with probability $p$, selects a random measurement basis, and measures the intercepted qubit by appending the appropriate gates and executing the circuit on the simulator backend. Eve then prepares a replacement qubit in the state corresponding to her measurement outcome and basis choice, including an identity gate channel marker for any noise on the Eve-to-Bob link, before forwarding it to Bob. All intercepted qubits are measured in a single batch call to the simulator for computational efficiency.

For BB84 and B92, Eve's random basis choice means she selects the correct basis with probability $\frac{1}{2}$. When she guesses incorrectly, her measurement collapses the qubit into a state incompatible with Alice's preparation, introducing a 50\% error probability on those rounds. The net effect for BB84 is a QBER contribution of $p \times 0.5 \times 0.5 = 0.25p$, reaching the characteristic 25\% at full interception ($p = 1$). For E91, Eve's interception degrades the entanglement fidelity of the Bell pairs, reducing the observed CHSH parameter $|S|$ towards the classical bound of 2. This provides an independent, physics-based eavesdropping detection mechanism that does not rely on QBER estimation.

The interception probability is swept from 0 to 1 in benchmarking experiments to generate the QBER-versus-interception-rate curves that form a central part of the results analysis.
% >>> END ADDED SECTION <<<


% >>> ADDED SECTION: Post-Processing Implementation <<<
\section{Post-Processing Implementation}

The post-processing pipeline is shared across all three protocols and consists of error estimation, secure key rate calculation, and mutual information computation.

\subsection{Error Estimation}

After sifting, the QBER is computed directly as the ratio of bit disagreements between Alice's and Bob's sifted keys to the total sifted key length. This computation is performed within the \texttt{QKDResult} dataclass, which stores both sifted keys and exposes the QBER as a derived property. The implementation does not explicitly sacrifice a subset of the sifted key for QBER estimation, as the simulation has access to both keys in their entirety; in a physical deployment, a configurable fraction would be publicly compared and discarded.

\subsection{GLLP Key Rate Calculation}

The secure key rate is computed using the GLLP formula presented in \cref{eq:gllp}. The binary Shannon entropy function is implemented within the \texttt{QKDResult} dataclass as a private method:
%
\begin{equation}
\label{eq:binary-entropy}
h(x) = -x \log_2(x) - (1 - x) \log_2(1 - x),
\end{equation}
%
with the convention that $0 \log_2(0) = 0$. The secure key rate property computes the secret fraction as $\max(0,\; 1 - h(\text{QBER}) - f_{EC} \cdot h(\text{QBER}))$ and multiplies it by the sifting efficiency (key rate) to yield the net secure key rate per transmitted qubit. The error correction efficiency is set to $f_{EC} = 1.16$ throughout, corresponding to a practical CASCADE implementation. This choice shifts the BB84 QBER security threshold from the idealised $\sim$11\% (at $f_{EC} = 1.0$) to $\sim$9.8\%, and reduces the B92 threshold to $\sim$6.5\%. These thresholds are the QBER values at which $R = 0$, beyond which no secret key can be extracted.

\subsection{Mutual Information}

For the mutual information analysis, the platform computes the Alice-Bob mutual information $I(A{:}B) = 1 - h(\text{QBER})$ and the Alice-Eve mutual information $I(A{:}E) = h(\text{QBER})$ as functions of the observed QBER. The crossover point where $I(A{:}E) \geq I(A{:}B)$ corresponds to the threshold beyond which secure key generation is impossible. This computation provides an alternative visualisation of the security threshold that complements the GLLP key rate analysis.
% >>> END ADDED SECTION <<<


% >>> ADDED SECTION: Performance Metrics <<<
\section{Performance Metrics}
\label{sec:methodology-metrics}

The benchmarking framework evaluates protocol performance using the following metrics, each of which is computed per simulation run and aggregated across multiple runs to produce confidence intervals.

The \textit{Quantum Bit Error Rate} (QBER) is defined as the ratio of bit errors in the sifted key to the total sifted key length. It is the primary indicator of channel quality and eavesdropper presence.

The \textit{sifting efficiency} is the ratio of sifted key length to total transmitted qubits. Theoretical targets are 50\% for BB84, 25\% for B92, and $2/9 \approx 22.2\%$ for E91 based on the sift table configuration.

The \textit{secure key rate} is the output of the GLLP formula (\cref{eq:gllp}), expressed as bits of secure key per sifted bit. A positive key rate indicates that secure key generation is feasible at the given noise level.

For E91, the \textit{CHSH parameter} $S$ is computed from the measurement correlations across non-matching-basis rounds, as defined in \cref{eq:chsh-computation}. A value of $|S| > 2$ confirms the presence of quantum correlations; degradation towards $|S| = 2$ indicates potential eavesdropping.

Each protocol class also provides a \texttt{theoretical\_qber()} method that returns the analytical QBER as a function of noise type and strength. For BB84, the theoretical QBER under depolarising noise is $p/2$. For B92, the depolarising theoretical QBER is $p/(1+p)$. These analytical predictions are overlaid on simulation data in the benchmarking plots as a validation mechanism.

\subsection{Statistical Methodology}

Each data point in the benchmarking results represents the mean of $N$ independent simulation runs (default: 30, configurable via the YAML configuration), each using $n$ qubits per run (default: 1000, configurable). The \texttt{BenchmarkRunner} collects the QBER, key rate, and secure key rate from each trial and computes the mean and standard deviation across trials. Error bars in the visualisations represent the standard error of the mean (SEM), computed as $\text{SEM} = \sigma / \sqrt{N}$, providing a measure of the uncertainty in the mean estimate rather than the spread of individual trials.

% NOTE: Update N and n values above to match your actual simulation parameters.
% >>> END ADDED SECTION <<<

